% year: תש"פ
% type: מבחן
% semester: ב
% moed: ב
% number: 1א
% points: 10
% categories: derivatives
% source: exams/מבחנים/תשפ/מועד ב סמסטר ב תשפ.pdf

\begin{question}
חשבו את הנגזרת של הפונקציה $f(x)=\sin(2x+1)$ בנקודה $x$ לפי הגדרת הנגזרת (תשובה לפי טבלת הנגזרות לא תתקבל).
\end{question}

\begin{hint}
השתמשו בזהות $\sin\alpha-\sin\beta=2\sin\frac{\alpha-\beta}{2}\cos\frac{\alpha+\beta}{2}$ ובגבול $\lim_{t\to0}\frac{\sin t}{t}=1$.
\end{hint}

\begin{solution}
לפי ההגדרה, $f'(x)=\lim_{h\to0}\frac{f(x+h)-f(x)}{h}$. לפי הזהות לסכום והפרש פונקציות,
\[ \sin(2x+2h+1)-\sin(2x+1)=2\sin(h)\cos(2x+1+h). \]
לכן
\[ f'(x)=\lim_{h\to0}\frac{2\sin h\cos(2x+1+h)}{h}=2\lim_{h\to0}\frac{\sin h}{h}\cdot\lim_{h\to0}\cos(2x+1+h)=2\cdot1\cdot\cos(2x+1), \]
כאשר השתמשנו בגבול היסודי $\frac{\sin h}{h}\to1$ וברציפות $\cos$. כלומר $\boxed{f'(x)=2\cos(2x+1)}$.
\end{solution}
